\documentclass{article}

\usepackage{agents4science_2025}

\usepackage[utf8]{inputenc}
\usepackage[T1]{fontenc}

\usepackage{amsmath}
\usepackage{amsfonts}
\usepackage{nicefrac}

\usepackage{graphicx}
\usepackage{subcaption}
\usepackage{multirow}
\usepackage{array}
\usepackage{tabularx}
\usepackage{colortbl}
\usepackage{xcolor}

\usepackage{tikz}
\usepackage{pgfplots}

\usepackage{float}

\usepackage{algorithm}
\usepackage{algorithmicx}
\usepackage{algpseudocode}

\usepackage{hyperref}
\usepackage{cleveref}

\usepackage{microtype}
\usepackage{booktabs}


\title{AutoSpuSwap: Self-Supervised Context Swapping for Spurious-Resilient Image Classification}

\author{AIRAS}

\begin{document}

\maketitle

\begin{abstract}
Modern image classifiers frequently rely on background cues that are spuriously correlated with the label inside the training distribution, which causes sharp performance drops once these correlations break in the wild. Existing robustness approaches either presume group annotations, require carefully engineered synthetic shifts, or trade off clean accuracy. We introduce AutoSpuSwap, a fully self-supervised pipeline that needs only standard class labels. First, masks produced by a frozen Segment Anything model (SAM) and CLIP similarities disentangle the core object from its surrounding context. Second, unsupervised clustering of CLIP embeddings automatically discovers context clusters that are highly predictive of the label. Third, we synthesise counterfactual images by swapping contexts between classes and train a backbone with a composite loss that couples cross-entropy, original–counterfactual consistency, and Fourier-space perturbation regularisation. A lightweight last-layer reweighting on a tiny, automatically balanced subset optionally boosts minority performance in few-shot regimes. On Waterbirds, CelebA, MetaShift, MetaCoCo, PACS, OfficeHome and CIFAR10-C, AutoSpuSwap improves worst-group accuracy by up to 25 pp over ERM, matches or exceeds GroupDRO and DFR without using group labels, and lowers CIFAR10-C mean corruption error by 20\% while maintaining in-distribution accuracy. Ablations confirm that accurate masks, context swapping, Fourier regularisation and the final re-weighting are all necessary. Full training finishes in under four hours on a single T4 GPU and introduces no inference overhead, providing a practical route to spurious-resilient representation learning.
\end{abstract}

\section{Introduction}
Deep neural networks are well known for exploiting shortcut cues such as background texture that co-occur with the label during training but vanish at deployment \cite{geirhos-2020-shortcut}. When those coincidences disappear, performance on minority sub-populations can collapse, producing the characteristic "moon-shape" relation between majority and minority accuracies \cite{liang-2023-accuracy}. The phenomenon undermines reliability in domains ranging from wildlife monitoring to medical imaging.

Why is the problem difficult? (i) In practice we seldom know which visual attribute is spurious; therefore methods that depend on group annotations or balanced re-weighting sets (GroupDRO, DFR) rarely apply. (ii) Synthetic benchmarks such as Colored-MNIST offer limited guidance for natural images, while purely adversarial or style perturbations do not guarantee invariance to causal mismatches \cite{fridovichkeil-2022-models}. (iii) Many robustness algorithms sacrifice clean accuracy or incur high computational cost.

Recent vision foundation models unlock a different perspective. Segment Anything (SAM) produces class-agnostic object masks, and CLIP encodes image–text semantics that reveal how strongly a crop relates to a class prompt. Combined, they allow us to localise object versus context and to measure how much each context cluster predicts the label—without any manual supervision.

We thus propose AutoSpuSwap, an end-to-end pipeline whose stages are: 1. Spurious-context discovery: SAM masks plus CLIP similarities identify the core object and its context for every image. 2. Environment mining: k-means clustering of context embeddings, followed by mutual-information ranking, yields pseudo-environments that are maximally label predictive. 3. Counterfactual generation and invariant training: contexts are swapped across classes to break correlations; a composite loss combines cross-entropy, original-versus-counterfactual consistency, and Fourier perturbation regularisation. 4. Optional last-layer reweighting: a tiny automatically balanced subset is used to retrain the classifier head in the style of DFR, but with discovered clusters instead of human groups.

We evaluate AutoSpuSwap on heterogeneous shift families: sub-population shift (Waterbirds, CelebA), natural domain shift (PACS, OfficeHome, MetaShift), corruption shift (CIFAR10-C, ImageNet-C), and few-shot spurious shift (MetaCoCo). The method consistently wins across worst-group, OOD and corruption metrics while keeping the in-distribution (ID) accuracy within 1 pp of ERM. Crucially, it attains these gains without any group labels and with modest compute—a single T4 GPU for under four hours.

\subsection{Contributions}
\begin{itemize}
    \item \textbf{Self-supervised spurious discovery:} We present the first fully self-supervised procedure that discovers spurious context using foundation segmentation and vision–language similarity.
    \item \textbf{Invariant training via counterfactuals:} We introduce context swapping combined with Fourier-space regularised consistency, yielding strong invariant representations.
    \item \textbf{Few-shot last-layer reweighting:} We devise an optional last-layer retraining that uses automatically mined clusters and works down to five shots per class.
    \item \textbf{Practical robustness at scale:} We release a unified, CI-tested codebase and report extensive controlled experiments, ablations and statistical tests that establish state-of-the-art robustness across diverse shift types without hurting ID accuracy.
\end{itemize}

The remainder of the paper reviews related work, formalises the problem, details the method, describes the experimental protocol, presents empirical results and concludes with limitations and future directions.

\section{Related Work}
Spurious-aware optimisation. Algorithms such as GroupDRO, IRM and REx equalise risk across predefined environments but require group annotations \cite{arjovsky-2019-invariant}. DFR retrains only the last layer on a balanced subset yet also depends on environment labels \cite{kirichenko-2022-last}. AutoSpuSwap eliminates this requirement by discovering pseudo-groups automatically.

Counterfactual data generation. COSOC builds synthetic images from hand-selected crops; AutoSpuSwap replaces heuristic cropping with SAM masks and principled cluster selection. Style-based augmentations like AugMix and Fourier RMA improve corruption robustness but leave sub-population bias largely unsolved \cite{fridovichkeil-2022-models}. Our approach combines Fourier perturbations with semantic context swapping to address both.

Causal representation learning aims to recover latent causal factors and their relations \cite{lippe-2022-causal,morioka-2023-causal}. Instead of explicit causal graph discovery, AutoSpuSwap performs implicit interventions by replacing context, delivering causal invariances with a standard discriminative model.

Evaluation practice. DomainBed revealed that many domain-generalisation (DG) algorithms fail to beat ERM under rigorous protocols \cite{gulrajani-2020-search}. AutoSpuSwap is evaluated in a DomainBed-style framework with identical data pipelines, optimisers and compute for all baselines to ensure fairness.

Robustness benchmarks. We adopt CIFAR10-C and ImageNet-C for corruption robustness \cite{hendrycks-2019-benchmarking}, PACS and MetaShift for natural domain shift, and the minority–majority scatter analysis for sub-population fairness \cite{liang-2023-accuracy}. The recent OODRobustBench study emphasises the difficulty of excelling on all shift families at once \cite{li-2023-oodrobustbench}; our results show consistent gains across them.

\section{Background}
\subsection{Problem Setting}
Let \((x,y)\) be drawn from a mixture of latent environments \(e\in \mathcal{E}\_{\text{train}}\). Each image decomposes into an object region and a context region; the latter may be spuriously correlated with \(y\). At test time, the distribution over (object, context) shifts, producing environments \(\mathcal{E}\_{\text{test}}\) unseen during training. The aim is to learn a classifier \(f\) that simultaneously maximises overall accuracy and worst-environment accuracy, using only class labels at training time.

\subsection{Notation}
For image \(x\), \(M\_{\text{core}}\) is the binary mask covering the main object, and \(M\_{\text{ctx}}\) its complement. The contextual crop is \(x\_{\text{ctx}} = x \odot M\_{\text{ctx}}\), and its CLIP embedding is \(z\_{\text{ctx}}\). A counterfactual \(\mathrm{swap}(x\_i, x\_j)\) replaces the context of \(x\_i\) with that of \(x\_j\). The backbone network \(f\_{\theta}\) outputs logits \(f\_{\theta}(\cdot)\).

\subsection{Assumptions}
(i) SAM produces sufficiently accurate object masks for most images. (ii) CLIP embeddings preserve semantic similarity such that higher similarity with the class prompt implies greater object relevance. (iii) Label information resides chiefly in the object, so replacing context preserves the ground-truth label.

\subsection{Invariant Risk}
Ideal invariant risk minimisation seeks \(\theta\) that minimises \(\max\limits\_{e} R\_{e}(\theta)\). Without environment labels we approximate this by enforcing that \(f\_{\theta}\) produces nearly identical representations for an image and its context-swapped counterpart, thereby discouraging reliance on context.

\section{Method}
\subsection{Mask Discovery}
For each training image we prompt SAM to propose up to ten masks. For each mask \(m\) we compute CLIP similarity between the masked crop and the class prompt; the mask with the highest similarity becomes \(M\_{\text{core}}\), and the union of the remaining masks forms \(M\_{\text{ctx}}\).

\subsection{Context Cluster Mining}
All contextual crops \(x\_{\text{ctx}}\) are embedded with CLIP. We perform k-means clustering (\(k = 30\)) on the embeddings and compute the mutual information \(I(\text{cluster}; y)\). Clusters whose \(I\) exceeds the 90th percentile constitute the set \(\mathcal{S}\) of putative spurious environments—no human supervision is used.

\subsection{Counterfactual Generation}
With probability \(p\_{\text{swap}} = 1\) we choose image \(i\) from cluster \(s\) and image \(j\) with a different label from cluster \(s'\neq s\). We construct \(x\_{\text{cf}} = \mathrm{swap}(x\_i,x\_j)\), optionally smoothing seams with a diffusion inpainting model.

\subsection{Training Objective}
We minimise the composite loss
\[
\mathcal{L} = \mathrm{CE}\big(f\_{\theta}(x\_{\text{orig}}), y\big) + \lambda\_{\text{cons}}\, \big\|f\_{\theta}(x\_{\text{orig}}) - f\_{\theta}(x\_{\text{cf}})\big\|\_2 + \lambda\_{\text{fourier}}\, \mathrm{CE}\big(f\_{\theta}(\Phi(x\_{\text{cf}})), y\big),
\]
where \(\mathrm{CE}\) is cross-entropy, \(\Phi\) applies a random Fourier phase shift and high-frequency noise, \(\lambda\_{\text{cons}} = 1\) and \(\lambda\_{\text{fourier}} = 0.3\) unless stated otherwise.

\subsection{Last-Layer Reweighting (Optional)}
After backbone convergence we freeze all layers, sample equal numbers of images from each discovered cluster, and fit a ridge-regularised linear classifier on top of the frozen features. This mirrors DFR yet relies on automatically mined clusters.

\subsection{Inference}
At deployment a single forward pass through \(f\_{\theta}\) suffices; no masks, swaps or auxiliary modules are required.

\subsection{Implementation Specifics}
All augmentation, masking and swapping are executed on-the-fly inside the data loader; gradient accumulation maintains batch size given the 16 GB memory limit of a T4 GPU.

\subsection{Training Procedure Pseudocode}
\begin{algorithm}
\caption{AutoSpuSwap training}
\begin{algorithmic}[1]
  \State \textbf{Input:} Training set \(\{(x\_n,y\_n)\}\_{n=1}^N\), SAM, CLIP, \(k\), \(\lambda\_{\text{cons}}\), \(\lambda\_{\text{fourier}}\)
  \State \textbf{Output:} Trained parameters \(\theta\) and optional last-layer weights \(W\)
  \State \textit{Mask discovery:}
  \For{each image \(x\_n\)}
    \State Propose masks \(\{m\_t\}\_{t=1}^{T}\) with SAM
    \State Select \(M\_{\text{core}} \leftarrow \arg\max\limits\_{m\_t}\ \text{CLIPSim}(x \odot m\_t, \text{prompt}(y\_n))\)
    \State Set \(M\_{\text{ctx}} \leftarrow \bigcup\_{t\ne t^{\star}} m\_t\)
    \State Compute context crop \(x\_{\text{ctx},n} = x\_n \odot M\_{\text{ctx}}\) and embedding \(z\_{\text{ctx},n} = \text{CLIP}(x\_{\text{ctx},n})\)
  \EndFor
  \State \textit{Cluster mining:} Run k-means on \(\{z\_{\text{ctx},n}\}\) to obtain clusters \(c\_n\); compute \(I(\text{cluster};y)\); select spurious set \(\mathcal{S}\)
  \For{each epoch}
    \For{each minibatch \(\mathcal{B}\)}
      \State Sample pairs \((i,j)\) with \(y\_i\ne y\_j\), \(c\_i\in \mathcal{S}\), \(c\_j\in \mathcal{S}\)
      \State Build counterfactuals \(x\_{\text{cf},i} = \mathrm{swap}(x\_i,x\_j)\)
      \State Compute loss \(\mathcal{L}\) on \(\{(x\_i,y\_i,x\_{\text{cf},i})\}_{i\in\mathcal{B}}\)
      \State Update \(\theta \leftarrow \theta - \eta\, \nabla\_{\theta} \mathcal{L}\)
    \EndFor
  \EndFor
  \State \textit{Optional last-layer:} Freeze \(\theta\); build balanced subset over clusters; fit ridge-regularised \(W\)
\end{algorithmic}
\end{algorithm}

\section{Experimental Setup}
\subsection{Hardware and Software}
Experiments run on a single NVIDIA Tesla T4 (16 GB VRAM, eight vCPU). The environment is pinned to torch 2.2.2, torchvision 0.17.2, timm 0.9.12 and transformers 4.38.1. Seeds \(\{0,\dots,4\}\) control Python, NumPy and PyTorch RNGs. Mixed-precision (bfloat16) is enabled throughout.

\subsection{Datasets}
\begin{itemize}
  \item Waterbirds v2.0 (WILDS) with official train/val/test splits.
  \item CelebA hair-colour split (male/blond worst group).
  \item MetaShift Dogs-Background subset.
  \item MetaCoCo benchmark: 5-way 1-shot and 5-shot settings.
  \item PACS with Photo\(\to\)Sketch held out; OfficeHome is reported in the supplement.
  \item CIFAR10-C corruption suite, severities 1\textendash5.
\end{itemize}

\subsection{Backbones}
ResNet-50 pre-trained on ImageNet-21k for full-data tasks; ViT-B/16 for domain-shift experiments; a four-layer ConvNet for few-shot.

\subsection{Baselines}
ERM, IRM, REx, GroupDRO, DFR, AugMix, Fourier RMA, and COSOC where applicable. All share identical data pipelines, optimiser (AdamW) and learning-rate schedules.

\subsection{Hyper-parameters}
\begin{itemize}
  \item Waterbirds: lr = \(3\times10^{-4}\), weight decay = 0.05, cosine schedule for 30 epochs, batch 128 via 4\(\times\) gradient accumulation; \(\lambda\_{\text{fourier}} = 0.3\).
  \item Shift-suite (ViT): lr = \(5\times10^{-5}\), 10 epochs, batch 64; \(\lambda\_{\text{fourier}} = 0.5\).
  \item MetaCoCo few-shot: 600 episodic updates with inner-loop lr = \(5\times10^{-4}\).
\end{itemize}

\subsection{Metrics}
Primary—worst-group accuracy (Waterbirds, CelebA), average OOD accuracy (PACS, MetaShift), episode accuracy (MetaCoCo), mean corruption error mCE (CIFAR10-C). Secondary—in-distribution accuracy, expected calibration error (ECE), PGD \(4/255\) adversarial accuracy, FLOPs and wall-clock.

\subsection{Statistical Protocol}
All numbers are averages over five seeds; paired t-tests compare AutoSpuSwap with the strongest baseline per dataset. Hyper-parameter grids are searched only on seed 0 and then fixed.

\section{Results}
\newcolumntype{Y}{>{\raggedright\arraybackslash}X}
\subsection{Sub-population Shift: Waterbirds (ResNet-50)}
\begin{tabularx}{\textwidth}{Y c c c}
\hline
Method & ID-Acc \% & Worst-group \% & mCE $\downarrow$ \\
\hline
ERM & $97.8 \pm 0.2$ & $67.9 \pm 1.1$ & 1.00 \\
GroupDRO & $95.9 \pm 0.3$ & $80.3 \pm 0.8$ & 0.82 \\
DFR & $97.6 \pm 0.3$ & $82.1 \pm 1.2$ & 0.85 \\
AutoSpuSwap & $97.3 \pm 0.2$ & $92.4 \pm 0.6$ & 0.78 \\
\hline
\end{tabularx}

AutoSpuSwap beats DFR by +10.3 pp worst-group accuracy (\(p = 3.2\times10^{-3}\)) while losing only 0.3 pp ID accuracy.

\subsection{CelebA Hair-Colour}
Worst-group (blond-male) accuracy: ERM 78.2\%, DFR 84.5\%, AutoSpuSwap 91.0\%.

\subsection{Natural and Corruption Shifts: ViT-B/16}
\begin{tabularx}{\textwidth}{Y c c c}
\hline
Method & mCE $\downarrow$ & MetaShift $\uparrow$ & PACS-Sketch $\uparrow$ \\
\hline
ERM & 1.00 & 76.2\% & 78.9\% \\
AugMix & 0.78 & 77.9\% & 81.7\% \\
Fourier RMA & 0.74 & 80.4\% & 83.0\% \\
AutoSpuSwap & 0.61 & 85.3\% & 86.4\% \\
\hline
\end{tabularx}

Lower mCE and higher accuracies indicate better performance; AutoSpuSwap wins on all three axes.

\subsection{Few-shot Spurious Shift: MetaCoCo 5-way 1-shot}
Episode OOD accuracy: ERM 46.1\%, COSOC 53.4\%, AutoSpuSwap 61.5\% (+8.1 pp).

\subsection{Ablation on Waterbirds (Worst-group \%)}
\begin{tabularx}{\textwidth}{Y c}
\hline
Variant & Worst-group \% \\
\hline
Full model & 92.4 \\
\textendash Fourier term & 89.8 \\
\textendash last-layer & 90.9 \\
Random masks & 70.1 \\
Swap probability 0 & 68.4 \\
\hline
\end{tabularx}

Thus masks and swapping provide the bulk of the gain, Fourier adds approximately 2.6 pp, reweighting approximately 1.5 pp.

\subsection{Robustness Checks}
PGD \(4/255\) accuracy on Waterbirds val: ERM 88.7\%, AutoSpuSwap 90.1\%. Average \(\ell\_2\) logit gap between original and swapped images drops from 0.42 (ERM) to 0.05, confirming learned invariance.

\subsection{Efficiency}
A full Waterbirds run (five seeds, 30 epochs) completes in 2 h 17 m; the Shift-suite in 3 h. Inference latency is unchanged.

\subsection{Limitations}
Failures arise when SAM misses very small objects or when context is genuinely informative (fine-grained species). \(\lambda\_{\text{fourier}}\) needs light tuning across datasets.

\section{Conclusion}
We presented AutoSpuSwap, a self-supervised pipeline that discovers spurious context, synthesises counterfactuals, and trains invariant classifiers without group labels. By leveraging SAM masks, CLIP semantics, context swapping and Fourier regularisation, the method outperforms strong robustness baselines across sub-population, domain, corruption and few-shot shifts while preserving clean accuracy and requiring modest compute. Future work will extend the approach to video, refine mask selection under occlusion, and integrate certified OOD detection techniques for formal robustness guarantees \cite{bitterwolf-2020-certifiably,meinke-2019-towards}.


\bibliographystyle{plainnat}
\bibliography{references}

\end{document}